% Standalone TeX research note. Intentionally not compiled.
% Bibliography is inline; there are no external input or bibliography files.
\documentclass[11pt]{article}
\usepackage[T1]{fontenc}
\usepackage[utf8]{inputenc}
\usepackage{lmodern}
\usepackage[margin=1in]{geometry}
\usepackage{amsmath,amssymb,amsthm,mathtools}
\usepackage[hidelinks]{hyperref}
\usepackage{microtype}
\newtheorem{theorem}{Theorem}[section]
\newtheorem{proposition}[theorem]{Proposition}
\newtheorem{corollary}[theorem]{Corollary}
\newtheorem{lemma}[theorem]{Lemma}
\theoremstyle{definition}
\newtheorem{definition}[theorem]{Definition}
\newtheorem{remark}[theorem]{Remark}
\title{Evidence Driven Revision of Characterization Languages}
\author{}
\date{4 November 2026}
\begin{document}
\maketitle

\begin{abstract}
A failed characterization can call for more capacity in the current language,
a different coordinate system, different transition or observation operations,
or a different task interface. These responses have different proof obligations.
This note connects a certified semantic-operation obstruction to a concrete
rational recurrence, separates exact and approximate weighted-automaton
expressivity, and proves what finite evidence can and cannot justify.
Under explicit degree, noise, and domain assumptions, rational identification
supplies a verified candidate. A true-error comparison against an entire
bounded linear class then licenses a dynamically realized, initially active
replacement gate with endpoint compression. The gate need not decrease loss
monotonically all the way to its endpoint. A least shift-stable localization
specifies the observable extension needed by the recurrence and distinguishes
that extension from general reciprocal closure. These results form a
conditional interface for a task-specific elaborator; they neither discover
the right language automatically nor impose one universal categorical
presentation of elaboration.
\end{abstract}

\section{The question and the dependency between the results}

The question is how evidence can justify changing the language in which a
phenomenon is characterized, its interpretation, and the operations that can
legitimately be performed, while retaining a checkable account of what the
previous characterization meant. Evidence can include observations, supplied
mathematical laws, proofs, and resource constraints. It is not restricted to
active experiments or numerical fitting.

The central running case is the positive-domain recurrence
$x_{n+1}=x_n/(1+x_n)$ with $x_0=1$. Its reciprocal coordinate has the affine
update $z_{n+1}=z_n+1$, but the readout becomes $1/z_n$. This is an exact
coordinate transport within a suitable rational language. The behavior
$1/(n+1)$ has infinite Hankel rank, so moving from finite weighted automata
with linear readouts to this rational language is instead a strict expressive
expansion. These two changes must not be conflated.

The arguments form the following conditional chain:
\[
 \begin{gathered}
 \text{declared task semantics and certified operation domains}\\
 \Downarrow\\
 \text{expressive obstruction and a verified alternative candidate}\\
 \Downarrow\\
 \text{validated evidence bounds and class-relative true-error comparison}\\
 \Downarrow\\
 \text{realizable gated replacement and exact endpoint removal}.
 \end{gathered}
\]
Every arrow has substantive premises. Finite observations alone cannot
refute the union of all finite WFA sizes, and the exact expressive separation
does not imply poor approximation at a nonzero absolute tolerance.
Candidate discovery, the choice of error metric and resources, and the
scientific warrant for a governing law remain additional questions.

All results below are mathematical proofs or explicitly conditional
constructions, not claims of formal verification. The individual ingredients have close prior art, discussed in
Section~\ref{sec:prior}; no novelty or priority claim is made. The trace
semantics used for the example does not account for first-person significance,
relevance, attention, or the choice of what counts as a phenomenon. It is a
bounded test of characterization revision, not a reduction of the broader
phenomenological programme to prediction.

\section{Task semantics and certified operations}
\label{sec:semantics}

Write $\mathbb N=\{0,1,2,\ldots\}$. Let $\mathcal C$ be a category of task contexts. A map $u:c'\to c$ specifies an
admissible restriction or substitution of context. Let
\[
  \mathcal B:\mathcal C^{\mathrm{op}}\longrightarrow\mathbf{Cat}
\]
be the chosen semantic interface. The objects of $\mathcal B(c)$ are the
behaviors, mathematical structures, or answers whose meaning the task requires
us to preserve. In the sequence example, $c=N$ is an observation horizon,
$\mathcal B(N)$ is the discrete category on $\mathbb R^{N+1}$, and
$u^*$ truncates a longer sequence. At an additional context $\infty$,
take $\mathcal B(\infty)$ to be the discrete category on $\mathbb R^{\mathbb N}$.
Other applications may use parameter
substitutions, domain restrictions, theorem contexts, or changes of available
questions. These alternatives require their own explicit $\mathcal B$.

A characterization $D$ has an indexed category of descriptions
$\mathcal M_D:\mathcal C^{\mathrm{op}}\to\mathbf{Cat}$ and an interpretation
\[
  b_D:\mathcal M_D\Longrightarrow\mathcal B.
\]
For readability, the framework uses strict indexed categories and strictly commuting
squares. Specified coherent isomorphisms give the usual pseudofunctor version.
The description category is not just a parameter space. Its objects may carry
syntax, types, laws, interpretations, partial-operation domains, and proof
certificates. Its arrows say which translations of these data are admissible.
Where only extensional results are required, discrete description categories
suffice and prevent stronger categorical claims being inferred accidentally.

Evidence is a context-indexed collection of accepted constraints. Write
$\mathcal E(c)\subseteq\mathcal B(c)$ for the full replete subcategory satisfying
them, with $u^*\mathcal E(c)\subseteq\mathcal E(c')$. The compatible-description
category is the inverse image
\[
  \mathcal R_D(E)(c)=b_D(c)^{-1}\mathcal E(c).
\]
Evidence provenance is additional data, not erased by this notation. A finite
measurement with a justified error interval, a hypothesis about a governing
law, and a checked proof under stated premises give different warrants.
For example, accepting an equation for every time index is stronger than
observing finitely many instances. The notation records the resulting
constraint; it does not manufacture its epistemic justification.

\paragraph{Verification languages.}
If the characterization also supplies an institution, its signature category,
sentence functor, model functor, and satisfaction relation must satisfy the
usual change-of-signature condition. With a signature map
$\sigma:\Sigma\to\Sigma'$, this condition has the orientation
\[
 M'\models_{\Sigma'}\operatorname{Sen}(\sigma)(\varphi)
 \quad\Longleftrightarrow\quad
 \operatorname{Mod}(\sigma)(M')\models_\Sigma\varphi.
\]
The model reduct is contravariant. A compiler from old programs to new programs
is instead often covariant. It must not be called an institution comorphism
merely because both constructions preserve something called meaning.
Institution theory already supplies the central abstract treatment of
satisfaction under vocabulary change \cite{cc-goguen-burstall,cc-goguen-rosu}.
The additional issue here is whether a proposed characterization change
actually has compatible descriptions and operation implementations.

A semantics-preserving compilation $T:D\to D'$ over the fixed interface is an
indexed functor $T:\mathcal M_D\to\mathcal M_{D'}$ with
$b_{D'}T=b_D$. It therefore restricts to
\[
 T:\mathcal R_D(E)\longrightarrow\mathcal R_{D'}(E).
\]
This last statement is a useful but elementary consequence of the equation;
it is not by itself a theory of learning. If the interface itself changes,
one must give a comparison functor between the two interfaces and state
which old assertions are preserved or forgotten. Equality on a small common
interface does not prove equality of all meanings in either language.

Available operations require more than names in a signature. An operation
has a domain of admissible inputs, a semantic effect, and a proof obligation
establishing applicability. At a fixed context write
\[
  o:U\longrightarrow B,\qquad U\subseteq B=\operatorname{Ob}\mathcal B(c).
\]
An implementation in $D$ sends every description $m$ supplied with a valid
certificate that $b_D(m)\in U$ to a description $O_D(m)$ and satisfies
$b_D(O_D(m))=o(b_D(m))$. The verification theory must explain how the
certificate is checked. Operations and their domains must commute with the
declared context restrictions, when both sides are defined.

Transporting an old operation along $T$ requires both a map on its
applicability certificates and a commuting semantic square. Requiring actual
syntactic equality $TO_D=O_{D'}T$ is stronger than necessary; a specified
behavior-preserving isomorphism or verified equivalence may suffice. That
choice must be fixed before claiming coherence of repeated translations.

\begin{proposition}[An operation has an expressive obstruction]
\label{prop:cc-operation-obstruction}
Let the semantic image at the context under consideration be
\[
 I_D=\{b_D(m):m\in\operatorname{Ob}\mathcal M_D(c)\}\subseteq B.
\]
Let $C_D$ be the set of pairs $(m,\pi)$ consisting of a description and an
accepted applicability certificate, and put
$J_D=\{b_D(m):(m,\pi)\in C_D\}\subseteq I_D\cap U$.
If $b\in J_D$ but $o(b)\notin I_D$, no implementation of $o$ in $D$
can be correct on every certified applicable input. In particular, neither
an optimizer nor a different parameterization of the same description class
can supply it. Conversely, if $o(J_D)\subseteq I_D$, a set-theoretic
implementation on $C_D$ exists after choosing representatives. This converse
makes no claim about computability, new proof production, functoriality, or
coherent reindexing.
\end{proposition}
\begin{proof}
Choose a certified pair $(m,\pi)$ with $b_D(m)=b$.
A correct output $O_D(m,\pi)$ would have denotation
$o(b)$ and hence put $o(b)$ in $I_D$, a contradiction. Conversely, choose
one description for each behavior in $o(J_D)$ and use that choice
after applying $o$ to the input's denotation. The construction is only a
set-theoretic one, which explains the stated limitations.
\end{proof}

This result turns the semantic diagram into a falsifiable requirement. A
purported transport fails if a required operation on a certified input leaves
the old semantic image. A semantically valid but uncertified applicability
condition does not suffice when the verification theory is incomplete.
Conversely, equal semantic images do not establish an effective or
canonical translation: different descriptions can implement a behavior in
incompatible ways, and checking equality may itself be undecidable.

\paragraph{Three different claims.}
A definitional change has specified inverse translations, with the appropriate
round-trip and obligation-preservation proofs. An expressive expansion has a
semantics-preserving inclusion and strictly more possible behaviors. A change
of task changes the interface or what evidence is accepted. These can occur
together, but one must not use a proof of one as a proof of another.
Equivalence of model categories alone is not the definition of a syntactic
definitional extension; the designated definitions and proof translations
matter.

\section{The operational interface and positive-domain transport}
\label{sec:operational}

Fix a finite nonempty alphabet $\Sigma$. A globally defined deterministic
real-valued presentation consists of a set $X$, an initial point $x_0\in X$,
maps $T_a:X\to X$ for $a\in\Sigma$, and an observation map $o:X\to\mathbb R$.
For a word $w=a_1\cdots a_k$, put
\[
 T_w=T_{a_k}\circ\cdots\circ T_{a_1},\qquad
 f(w)=o(T_wx_0).
\]
The empty-word transition is the identity. The word ``presentation'' here
does not prescribe a universal syntax or category of presentations. It
records only the semantic interface used by the following theorems.
Different typed, contextual, dependent, or diagrammatic elaboration
infrastructures can discharge the required existence and preservation
obligations in different ways.

Real weighted finite automata (WFAs) are one presentation language: use
row-vector states, $x_0=\alpha^{\mathsf T}$, $T_a(x)=xA_a$, and
$o(x)=x\beta$. Writing $A_w=A_{a_1}\cdots A_{a_k}$ for
$w=a_1\cdots a_k$ gives $f(w)=\alpha^{\mathsf T}A_w\beta$. Rational state updates with a specified
forward-invariant domain form another. A mathematical formula with a pole
on a reachable state is \emph{not} a globally defined presentation.
A trajectory-certified rational description can be interpreted on its
forward-invariant reachable state set. A language requiring a finitely
presented invariant domain has a further obligation to construct and verify
that domain; the extensional presentation alone does not discharge it.

For finite prefix and suffix sets $S,E\subseteq\Sigma^*$, write
\[
 H_f[S,E]_{s,e}=f(se).
\]
For a WFA of dimension $r$, this block factors as
\[
 H_f[S,E]=P Q,\qquad
 P_{s,:}=\alpha^{\mathsf T}A_s,\quad Q_{:,e}=A_e\beta,
\]
where $P$ has $r$ columns and $Q$ has $r$ rows.
Thus every such block has rank at most $r$.
Every block comparison below uses the same $S,E$ and scalar observations
for the two languages being compared.
Repeated occurrences of a word must have the same predicted value.
An alternative language that changes the observations or interventions
needs an explicit comparison map; its raw loss is not automatically
comparable with the one used here.

\begin{remark}[Two distinct meanings of a global assertion]
A candidate presentation can define $f(w)$ for \emph{all} words, and the
coexistence construction below preserves its semantic equations for all
words. Neither fact proves that the candidate agrees with an unknown
target on all words. A finite-block comparison is a statement about that
block. Extending its statistical or scientific interpretation is a
separate obligation.
\end{remark}

Consider the recurrence already motivating the comparison:
\begin{equation}
\label{eq:cc-rational}
 x_0=1,\qquad x_{n+1}=\frac{x_n}{1+x_n},\qquad y_n=x_n.
\end{equation}
Its state domain is $(0,\infty)$. The transition maps this domain into itself,
so every division is defined. The reciprocal coordinate $z=1/x$ is a bijection
of the positive domain with itself, with the same formula for its inverse.
It yields
\begin{equation}
\label{eq:cc-affine}
 z_0=1,\qquad z_{n+1}=z_n+1,\qquad y_n=\frac1{z_n}.
\end{equation}
In particular $z_n=n+1$ and $y_n=1/(n+1)$ by induction.

\begin{proposition}[Conjugacy transports the readout and operations]
\label{prop:cc-conjugacy}
Let $X,Z$ be state domains, $\phi:X\to Z$ a bijection, $F:X\to X$ a
transition, and $q:X\to Y$ a readout. Set
\[
 F'=\phi F\phi^{-1},\qquad q'=q\phi^{-1},\qquad z_0=\phi(x_0).
\]
Then $q'(F'^n(z_0))=q(F^n(x_0))$ for all $n\ge0$. Every partial state
operation $g:P\to X$ transports to
$g'=\phi g\phi^{-1}:\phi(P)\to Z$. Conjugacy preserves identities,
defined composition, and equations of operations; readout predicates
transport by inverse image. These identities commute with restriction to
every finite horizon.
\end{proposition}
\begin{proof}
Induction gives $F'^n=\phi F^n\phi^{-1}$, after which the readout identity
follows by cancellation. The corresponding cancellation proves the
composition identities for partial operations, with transported domains.
A predicate $q(x)\in V$ becomes $q'(z)\in V$, so truth is unchanged.
Horizon restriction discards some already equal readout values.
\end{proof}

For \eqref{eq:cc-rational}, the old observable $x$ must become $1/z$.
Leaving the readout equal to $z$ would produce $n+1$, a different sequence.
The old operation $x\mapsto x/(1+x)$ becomes $z\mapsto z+1$, and its positive
domain proof transports. More general old observables $p(x)$ become
$p(1/z)$. These may be rational even when $p$ was polynomial. Thus the
transformation is a definitional equivalence in a language allowing the
displayed rational operations and their domain predicates, but not an
equivalence inside a language restricted to linear readouts. The claimed
equivalence depends on the expressive profile of both transitions and
observations, not only on the state dimension.

There are two distinct uses of reciprocal here. A reciprocal \emph{state
coordinate} with a transported readout preserves the phenomenon. Taking
the reciprocal of an \emph{output sequence} deliberately changes the
phenomenon. The next obstruction concerns the second operation.

\begin{proposition}[A parameterized rational presentation]
\label{lr:rational-family}
For $c\geq0$ and $x_0>0$, the map
\[
 R_c(x)=\frac{x}{1+cx},\qquad x\in(0,\infty),
\]
with identity readout defines the behavior
\[
 b_c(a^n)=\frac{x_0}{1+ncx_0}.
\]
Reciprocal coordinates transport it to $z\mapsto z+c$, with initial state
$z_0=1/x_0$ and readout $z\mapsto1/z$.
\end{proposition}
\begin{proof}
The denominator is positive on the stated domain, and $R_c$ preserves
positivity. The identity $1/R_c(x)=1/x+c$ proves the transport and the
behavior formula by induction.
\end{proof}

This transport preserves the observed behavior. It does not identify the
underlying physical mechanism. In particular, the transition being affine
in the new coordinate does not make the complete presentation a linear WFA:
its reciprocal readout must also be charged to the language.

\begin{proposition}[The revised language generates composition laws and tests]
\label{lr:generated-tests}
For the family $R_c$ with $c\geq0$ on $(0,\infty)$,
\[
 R_0=\operatorname{id},\qquad R_c\circ R_d=R_{c+d}.
\]
Suppose each actually available action $a\in\Sigma$ is hypothesized
to act as $R_{c_a}$ for some $c_a\geq0$, with the same state coordinate
and observation $o(x)=x$. The hypothesis predicts
\[
 b(w)=\frac{x_0}{1+x_0\sum_{i=1}^{|w|}c_{a_i}}
 \quad(w=a_1\cdots a_{|w|}),
 \qquad b(hab)=b(hba)
\]
for every history $h$ and every pair of actions $a,b$.
If the two word experiments $hab,hba$ are executable from the same
initial preparation, and their measured outputs have absolute error
at most $\eta$, an observed difference greater than $2\eta$ refutes
this \emph{entire common-coordinate rational-action family}, regardless
of the fitted parameter values.
\end{proposition}
\begin{proof}
Direct substitution gives
$R_c(R_d(x))=x/(1+(c+d)x)$. Equivalently, reciprocal coordinates
transport composition to addition of translation increments. Iterating
proves the word formula; addition is commutative. Under the hypothesis
the two true outputs coincide, so their measured difference is at most
the sum of their two error bounds.
\end{proof}

The point is the direction of dependence: a proposed language yields an
operation law, that law yields a test of the whole language, and a failed
test can send the learner back to representation revision. This is not
merely another parameter fit. It does not assert that experiment design
is solved: other alternatives may satisfy the same identity, and the
requisite actions and common preparation must actually be available.
Nor does the equality claim follow from finitely many arbitrary
observations without the rational-family hypothesis. It is a derived
obligation of that hypothesis, suitable for the verification and
elaboration infrastructure chosen for the task.

\section{Exact and approximate expressive separation}
\label{sec:expressivity}

A real one-letter weighted finite automaton of dimension $r$ computes
\[
 f_n=\alpha^{\mathsf T}A^n\beta,
 \quad A\in\mathbb R^{r\times r},\quad \alpha,\beta\in\mathbb R^r.
\]
Write $\mathrm{WFA}_{\mathrm{fin}}$ for the union over finite $r$.
A rational state-space language permits a finite-dimensional state domain,
rational transition and readout expressions, and certificates that the
trajectory remains where all denominators are nonzero. This language
contains the WFA descriptions by taking row states, $F(v)=vA$, $v_0=\alpha^{\mathsf T}$, and
$q(v)=v\beta$. This inclusion preserves every sequence value,
the shift operation, and finite-horizon restriction. Its existence does not
provide a converse compilation. For $m\geq1$, use the consistent
square-block notation
\[
 H_m=\left(\frac1{i+j+1}\right)_{0\leq i,j<m}.
\]

\begin{theorem}[No finite WFA realizes the reciprocal trajectory]
\label{thm:cc-hilbert}
The sequence $f_n=1/(n+1)$ has no finite-dimensional real WFA realization.
For every $r$, its values $f_0,\ldots,f_{2r}$ already exclude all WFA
realizations of dimension at most $r$.
\end{theorem}
\begin{proof}
For any $r$-state WFA, the $(r+1)\times(r+1)$ Hankel block factors as
\[
 H_{ij}=f_{i+j}
  =(\alpha^{\mathsf T}A^i)(A^j\beta),\quad 0\le i,j\le r,
\]
so its rank is at most $r$. For the specified sequence and any nonzero
$v\in\mathbb R^{r+1}$,
\[
 v^{\mathsf T}Hv
 =\sum_{i,j=0}^r\frac{v_iv_j}{i+j+1}
 =\int_0^1\left(\sum_{i=0}^r v_i t^i\right)^2dt>0.
\]
A nonzero polynomial is nonzero on some open subinterval, proving strict
positivity. Thus $H$ has rank $r+1$, contradicting any such realization.
Since $r$ was arbitrary, no finite WFA realizes the whole sequence.
\end{proof}

This is an elementary instance of the classical connection between WFA
realizations and finite Hankel rank; it is not a new realization theorem
\cite{cc-balle-mohri}. With an observed block $\widehat H$ and a justified
operator-norm error bound $\|H-\widehat H\|_2\le\epsilon$, the stronger
condition $\sigma_{r+1}(\widehat H)>\epsilon$ gives a noise-robust exclusion.
The exact integral proof does not supply that measurement bound.

\begin{corollary}[Reciprocal is unavailable in the finite WFA language]
\label{cor:cc-reciprocal}
The class $\mathrm{WFA}_{\mathrm{fin}}$ is not closed under pointwise
reciprocal of nowhere-zero sequences. The rational state-space language
implements this operation whenever its nonzero-output obligation is proved.
Hence a semantics-preserving implementation of this operation cannot be
transported back to all finite WFA descriptions, even with unbounded but
finite dimension chosen separately for each input.
\end{corollary}
\begin{proof}
The sequence $g_n=n+1$ is realized by
\[
 A=\begin{pmatrix}1&1\\0&1\end{pmatrix},\quad
 \beta=\begin{pmatrix}1\\1\end{pmatrix},\quad
 \alpha=\begin{pmatrix}1\\0\end{pmatrix}.
\]
The identity $g_n=n+1$ and its positivity have elementary induction proofs,
taken here as accepted applicability certificates for this example.
Its pointwise reciprocal is the sequence excluded by
Theorem~\ref{thm:cc-hilbert}. Proposition~\ref{prop:cc-operation-obstruction}
therefore applies. In a rational realization with readout $q$, replace the
readout by $1/q$ and retain the state dynamics. The operation is defined
on the reachable trajectory exactly when the stated obligation holds.
\end{proof}

Consequently the WFA-to-rational inclusion is a genuine expressive expansion.
The $x$-to-$z$ transformation within the rational language is a definitional
change. Conflating these two arrows would hide precisely the distinction
that a theory of characterization change needs to expose.

\begin{proposition}[Uniform absolute error bounds]
\label{prop:cc-approximation}
For every $k\ge1$, the $k$-state diagonal WFA
\[
 g_n^{(k)}=\frac1k\sum_{j=1}^k\left(\frac jk\right)^n
\]
satisfies $\sup_{n\ge0}|g_n^{(k)}-1/(n+1)|\le1/k$.
Conversely, for every $r$-state WFA sequence $g$,
\[
 \sup_{n\ge0}|g_n-1/(n+1)|
 \ \ge\ \frac{\sigma_{r+1}(H_{r+1})}{r+1}>0,
 \qquad (H_{r+1})_{ij}=\frac1{i+j+1}\quad(0\le i,j\le r).
\]
\end{proposition}
\begin{proof}
For $n\ge1$, the function $t^n$ is increasing on $[0,1]$. Its right and
left Riemann sums on the uniform $k$-interval partition bracket its integral
and differ by $1/k$. The right sum is $g_n^{(k)}$, while the integral is
$1/(n+1)$. At $n=0$, both the target and the approximation equal one.
The diagonal realization uses eigenvalues $j/k$ and weights $1/k$.

For the lower bound, the corresponding $(r+1)$-square block $G_{r+1}$ of $g$
has rank at most $r$. Choose a unit vector $v\in\ker G_{r+1}$. Then
$\|H_{r+1}-G_{r+1}\|_2\ge\|H_{r+1}v\|_2\ge\sigma_{r+1}(H_{r+1})$. If the uniform
entrywise error is $\epsilon$, the block's Frobenius norm and hence its
operator norm are at most $(r+1)\epsilon$. Combine these inequalities and
use the positive definiteness proved above.
\end{proof}

The upper construction is elementary rather than optimal. Braess and
Hackbusch prove uniform approximation of $1/x$ on $[1,\infty)$ by
$k$-term exponential sums with error $O(\exp(-c\sqrt{k}))$ for a positive
constant $c$ \cite{cc-braess-hackbusch}. Sampling a sum
$\sum_j w_j e^{-a_jx}$ at $x=n+1$ gives a $k$-state diagonal WFA with
eigenvalues $e^{-a_j}$ and weights $w_je^{-a_j}$, so their bound applies
to this sequence as well. Their conjectured optimal constant in the exponent
is not treated as a proved bound.

Thus the exact infinite-rank obstruction is compatible with useful
finite-state approximation, even uniformly over all time indices in
absolute error. The fixed-rank lower bound above can be very small because
Hilbert matrices are poorly conditioned. A practical choice of language
must therefore specify accuracy, metric, resource cost, and which symbolic
operations or explanations the task requires. A claim about exact equality
does not establish practical inferiority at a chosen nonzero tolerance;
uniform absolute and uniform relative error are also different requirements.

\begin{proposition}[Global relative error gives a different obstruction]
\label{prop:relative-obstruction}
For every finite real WFA sequence $g$,
\[
 \sup_{n\geq0}\frac{|g_n-1/(n+1)|}{1/(n+1)}\geq1.
\]
Thus no finite WFA achieves a uniform relative error strictly below one,
even though finite WFAs achieve arbitrarily small uniform absolute error.
The lower bound is sharp: the zero behavior has relative error exactly one.
\end{proposition}
\begin{proof}
Suppose instead that $|g_n-1/(n+1)|\leq\varepsilon/(n+1)$ for every
$n$, where $0\leq\varepsilon<1$. The power series
$G(t)=\sum_{n\geq0}g_nt^n$ converges absolutely for $|t|<1$, and for
$0<t<1$ obeys
\[
 (1-\varepsilon)\frac{-\log(1-t)}t
 \leq G(t)\leq
 (1+\varepsilon)\frac{-\log(1-t)}t.
\]
A finite WFA has the rational formal generating function
$\alpha^{\mathsf T}(I-tA)^{-1}\beta$. The power-series identity near
zero, followed by analytic continuation, identifies its reduced rational
function with $G$ on the unit disk; any apparent poles there are removable.
At $t=1$, a rational function is either bounded or has a pole of some
positive integer order. Boundedness contradicts the lower logarithmic
bound, while a pole cannot satisfy the two logarithmic bounds. This rules
out every such $\varepsilon$ and proves the assertion.
\end{proof}

The absolute/relative distinction is a change in the comparison interface,
not a contradiction between approximation results. A claim that one
language is adequate must state which of these requirements the task uses.

\section{Finite evidence and the inverse-limit obstruction}
\label{sec:finite}

\begin{proposition}[Finite non-refutation and an unseen-test obstruction]
\label{lr:finite-obstruction}
Let $D\subset\Sigma^*$ be finite.
\begin{enumerate}
\item Every assignment $y:D\to\mathbb R$ is the restriction of a
      finite real WFA behavior.
\item For any behavior $b$, any $v\notin D$, and any $M>0$, the
      behaviors $b$ and $b+M\mathbf1_{\{v\}}$ agree on $D$ but differ
      by $M$ at $v$. Any common prediction at $v$ has absolute error at
      least $M/2$ for one of them.
\item If $b$ has a finite WFA realization, so does the modified behavior.
      Thus the unseen-test obstruction persists even within the union
      of all finite WFA sizes.
\end{enumerate}
\end{proposition}
\begin{proof}
Extend the finite assignment by zero outside $D$. A trie automaton with
state set $\{\varepsilon\}\cup\operatorname{Pref}(D)$, where
$\operatorname{Pref}(D)$ contains the prefixes of words in $D$, initial state the empty
prefix, transitions following trie edges and otherwise to zero, and
terminal readout $y(w)$ on $w\in D$, realizes the extension. The same
construction realizes $\mathbf1_{\{v\}}$. Direct sums realize sums of
WFA behaviors. The prediction lower bound follows from the triangle
inequality: two real values distance $M$ apart cannot both be within
less than $M/2$ of the same prediction.
\end{proof}

The argument also applies to finitely many adaptively selected exact
queries: choose an unqueried word after the finite transcript, and the
two completions are indistinguishable on that transcript. Consequently,
finite evidence can refute a \emph{fixed} linear dimension budget without
refuting the entire language of finite WFAs. It can certify a candidate's
advantage on a specified block without identifying its unseen behavior
against unrestricted alternatives.

\begin{corollary}[Semantic image does not commute with the horizon limit]
\label{thm:cc-finite-prefix}
Let the description set at every horizon be the same set of complete
finite unary WFA descriptions, with identity restriction maps between
description sets. Interpret a description at horizon $N$ by its prefix of
length $N+1$, and at infinity by its complete sequence. If $I_N$ and
$I_\infty$ denote these semantic images, then
\[
 I_N=\mathbb R^{N+1},\qquad
 I_\infty\subsetneq\varprojlim_N I_N=\mathbb R^{\mathbb N}.
\]
\end{corollary}
\begin{proof}
Proposition~\ref{lr:finite-obstruction} realizes every finite prefix,
so $I_N=\mathbb R^{N+1}$. The inverse limit under truncation is the set
of all infinite sequences. Its element $1/(n+1)$ is absent from
$I_\infty$ by Theorem~\ref{thm:cc-hilbert}. A compatible family in the
constant description system, in contrast, is one fixed finite WFA.
Thus taking the semantic image after the limit of descriptions differs
from taking the limit of the finite semantic images.
\end{proof}

An explicit interpolant for $a_0,\ldots,a_N$ uses a shift register of
dimension $N+1$: in column notation take $Ae_i=e_{i+1}$ for $i<N$,
$Ae_N=0$, initial vector $e_0$, and readout $e_i\mapsto a_i$.
This realizes the specified prefix and is zero thereafter. Transposition
converts it to the row-state convention used above. The obvious inclusion
into the next register fails to intertwine the last transition, while the
obvious truncation projection fails to preserve the full readout when
$a_{N+1}\ne0$, as in the harmonic example. Compatible finite behaviors
therefore do not automatically supply compatible full model transports.

These obstructions do not prevent preferring a rational characterization.
Rejecting every finite WFA from finitely many values requires
something besides those values: a dimension or resource constraint, a supplied
law with an accepted universal scope, or an explicitly stated inductive
criterion. Four observed values do not turn a rational recurrence into an
established law of the target.

If the accepted evidence instead includes the finite \emph{presentation}
\[
 f_0=1,\qquad
 \forall n\in\mathbb N:\ f_n>0\quad\text{and}\quad
 f_{n+1}(1+f_n)=f_n,
\]
ordinary induction identifies its unique sequence as $1/(n+1)$. The
presentation is finite syntax with infinite scope; it is not a finite table
of measurements. Under that premise, the WFA characterization has an empty
compatible fiber, while the rational characterization has the verified
realization \eqref{eq:cc-rational}. This is a solved conditional problem.
Whether observed data justify accepting the premise remains a different,
statistical or scientific inference problem.

\section{Assumption-relative rational identification}
\label{sec:identification}

The input coordinate and its interpretation are fixed in this section.
For the running recurrence with identity readout, consecutive outputs supply
transition pairs. With an unknown hidden-state coordinate or unknown
readout, identifying the following rational map from given pairs does not
solve joint state and observation identification.

Let $\mathcal R_{p,q}$ be the rational functions representable as $P/Q$
with $\deg P\leq p$, $\deg Q\leq q$, $Q\not\equiv0$. Evaluations are
only made where the denominator is nonzero.

\begin{theorem}[Finite uniqueness of a bounded-degree rational function]
\label{lr:rational-unique}
If $R,\widetilde R\in\mathcal R_{p,q}$ agree on $p+q+1$ distinct
inputs at which both are defined, then they are the same rational
function. This is a uniqueness statement, not an existence statement.
\end{theorem}
\begin{proof}
Write $R=P/Q$ and $\widetilde R=\widetilde P/\widetilde Q$.
The polynomial $P\widetilde Q-\widetilde P Q$ has degree at most
$p+q$ and has $p+q+1$ distinct roots. It is identically zero.
\end{proof}

For example the three pairs
\[
 (1,1/2),\qquad(1/2,1/3),\qquad(1/3,1/4)
\]
uniquely identify $x/(1+x)$ \emph{among degree-$(1,1)$ rational
functions}. The existence of an unconstrained continuation outside that
language is not excluded. In particular, this inference must not be
described as recovering an assumption-free law from three observations.

\paragraph{Coefficient charts and a hidden degeneracy.}
For numerical recovery, choose a chart with $Q(0)\ne0$ and normalize
\[
 P_\theta(x)=\sum_{j=0}^{p}a_jx^j,
 \quad Q_\theta(x)=1+\sum_{k=1}^{q}c_kx^k,
 \quad\theta=(a_0,\ldots,a_p,c_1,\ldots,c_q)^T.
\]
This normalization does not cover rational functions whose reduced
denominator vanishes at zero; another chart is needed there.
For exact pairs $(x_i,y_i)$, put
\[
 A_{i,:}=(1,x_i,\ldots,x_i^p,-y_ix_i,\ldots,-y_ix_i^q),
 \qquad y=(y_i)_{i=1}^{N}.
\]
The coefficient equations are $A\theta=y$. Function uniqueness does
not always imply coefficient uniqueness: common factors or slack degree
bounds can give different coefficient vectors for the same rational
function.

\begin{proposition}[An identifiable normalized coefficient chart]
\label{lr:coefficient-identifiable}
Suppose $P,Q$ are coprime, $Q(0)=1$, $\deg P=p$, and $\deg Q=q$.
At $N\geq p+q+1$ distinct inputs away from the poles, with
$y_i=P(x_i)/Q(x_i)$, the matrix $A$ has full column rank.
\end{proposition}
\begin{proof}
A null vector gives polynomials $\delta P,\delta Q$ of degrees at most
$p,q$, with $\delta Q(0)=0$, and
$\delta P(x_i)Q(x_i)-P(x_i)\delta Q(x_i)=0$.
The polynomial identity argument makes this equality global.
Coprimality implies $Q$ divides $\delta Q$. Since $\deg Q=q$,
$\delta Q=tQ$ for a constant $t$. Its zero constant coefficient gives
$t=0$, and then $\delta P=0$.
\end{proof}

The degree assumptions are sufficient, not necessary. Rank should be
checked in the actual chart. If a larger degree bound is used than the
data-generating rational function needs, redundant common-factor
directions may destroy coefficient conditioning without destroying
uniqueness of the represented function.

\begin{theorem}[Conditioned noisy interpolation]
\label{lr:noisy-interpolation}
Assume exact data satisfy $A\theta_*=y$, with
$\|\theta_*\|_2\leq M$. Let measured pairs produce
$\widehat A=A+E_A$ and $\widehat y=y+e_y$, with justified bounds
\[
 \|E_A\|_2\leq\delta_A,
 \qquad\|e_y\|_2\leq\delta_y.
\]
If $\widehat A$ has full column rank and
$\widehat\theta=\widehat A^\dagger\widehat y$, then
\begin{equation}
 \|\widehat\theta-\theta_*\|_2
 \leq\kappa:={\delta_y+\delta_A M\over
                     \sigma_{\min}(\widehat A)}.
 \label{lr:kappa}
\end{equation}
If additionally $\sigma_{\min}(\widehat A)>\delta_A$, the exact
matrix $A$ is also certified to have full column rank.
\end{theorem}
\begin{proof}
Full column rank gives $\widehat A^\dagger\widehat A=I$. Therefore
\[
 \widehat\theta-\theta_*
 =\widehat A^\dagger(e_y-E_A\theta_*).
\]
Taking norms proves \eqref{lr:kappa}. The final assertion follows from
singular-value perturbation.
\end{proof}

The error bounds can include noise in \emph{both} coordinates of each
pair. For example suppose $|x_i|,|\widehat x_i|\leq R$ with $R\geq1$,
$|y_i|\leq B_y$, $|\widehat x_i-x_i|\leq\eta_x$, and
$|\widehat y_i-y_i|\leq\eta_y$. Then valid choices are
\begin{align}
 \delta_y&=\sqrt N\eta_y,\\
 \delta_A&=\sqrt N\left[
 \sum_{j=1}^{p}(jR^{j-1}\eta_x)^2+
 \sum_{k=1}^{q}(R^k\eta_y+B_y kR^{k-1}\eta_x)^2
 \right]^{1/2}.
 \label{lr:design-noise}
\end{align}
Indeed the mean-value theorem bounds monomial perturbations, and
\[
 |\widehat y_i\widehat x_i^k-y_ix_i^k|
 \leq R^k\eta_y+B_y kR^{k-1}\eta_x.
\]
The resulting Frobenius bound controls the operator norm. Correlations
between consecutive noisy transition pairs do not affect this
deterministic calculation; they do matter when establishing the noise
event probabilistically.

\begin{proposition}[From coefficients to a safe rational function]
\label{lr:safe-domain}
Let $I\subseteq[-R,R]$, and suppose on $I$ the true normalized
rational function obeys
\[
 |Q_{\theta_*}(x)|\geq\gamma>0,
 \qquad |R_{\theta_*}(x)|\leq B.
\]
Put
\[
 V_p(R)=\left(\sum_{j=0}^{p}R^{2j}\right)^{1/2},
 \qquad W_q(R)=\left(\sum_{k=1}^{q}R^{2k}\right)^{1/2},
\]
with $W_0=0$. If $\|\widehat\theta-\theta_*\|_2\leq\kappa$
and $\kappa W_q(R)<\gamma$, then $\widehat R$ has no pole on $I$
and
\begin{equation}
 \sup_{x\in I}|\widehat R(x)-R_{\theta_*}(x)|
 \leq\zeta:={\kappa\sqrt{V_p(R)^2+B^2W_q(R)^2}
                       \over\gamma-\kappa W_q(R)}.
 \label{lr:rational-error}
\end{equation}
\end{proposition}
\begin{proof}
The denominator perturbation is at most $\kappa W_q(R)$.
Writing $\delta P=\widehat P-P$ and $\delta Q=\widehat Q-Q$ gives
\[
 \widehat R(x)-R_{\theta_*}(x)
 ={\delta P(x)-R_{\theta_*}(x)\delta Q(x)\over\widehat Q(x)}.
\]
The numerator is the inner product of the coefficient error with
$(1,x,\ldots,x^p,-R_{\theta_*}(x)x,\ldots,-R_{\theta_*}(x)x^q)$,
whose norm is bounded by the
square root in \eqref{lr:rational-error}.
\end{proof}

No-pole control alone is not forward invariance. One must separately
verify $\widehat R(I)\subseteq I$, enforce it in a candidate family, or
prove that the particular reachable orbit stays in a safe domain.
For an interval $I$, a sufficient condition is that $R_{\theta_*}(I)$ has distance
greater than $\zeta$ from the complement of $I$. That sufficient condition
does not apply when the true dynamics touch the boundary. These are real
verification obligations, not consequences of a low interpolation loss.

\begin{proposition}[Conditional rollout control]
\label{lr:rollout}
Suppose $R,\widehat R:I\to I$, $R$ is $K$-Lipschitz on $I$, and
$\sup_I|\widehat R-R|\leq\zeta$. For their trajectories from
$x_0,\widehat x_0\in I$,
\[
 |\widehat x_n-x_n|
 \leq K^n|\widehat x_0-x_0|+
       \zeta\sum_{j=0}^{n-1}K^j.
\]
If instead $R=R_c$, $\widehat R=R_{\widehat c}$ have the same
$x_0>0$ and $c,\widehat c\geq c_{\min}>0$, then
\begin{equation}
 \sup_{n\geq0}|b_c(a^n)-b_{\widehat c}(a^n)|
 \leq {x_0\over4c_{\min}}|c-\widehat c|.
 \label{lr:uniform-special-family}
\end{equation}
\end{proposition}
\begin{proof}
Add and subtract $R(\widehat x_n)$ to obtain the one-step error
recursion $e_{n+1}\leq Ke_n+\zeta$, then iterate it. For the special
family the exact difference is
\[
 {n x_0^2|c-\widehat c|\over
 (1+ncx_0)(1+n\widehat c x_0)}
 \leq {x_0|c-\widehat c|\over c_{\min}}
       {t\over(1+t)^2},\qquad t=nc_{\min}x_0.
\]
Use $t/(1+t)^2\leq1/4$.
\end{proof}

The uniform estimate is stronger than generic finite-horizon control
because it uses the special rational family. Within that narrower family,
one exact nonzero transition pair determines
$c=1/y-1/x$; this statement carries a correspondingly stronger language
assumption than degree-$(1,1)$ rational interpolation. If both observed
and true $x,y$ are bounded below by $a>0$, coordinatewise errors at most
$\eta$ give $|\widehat c-c|\leq2\eta/a^2$. Positivity of
$\widehat c$ must still be checked.

\section{From a true-error certificate to a realized language change}
\label{sec:gate}

Let $H_*\in\mathbb R^{m\times n}$ be a true observation block and let
$Y$ be its measured version. Assume justified bounds
\begin{equation}
 \|Y-H_*\|_2\leq\epsilon_2,
 \qquad \|Y-H_*\|_F\leq\epsilon_F.
 \label{lr:noise-event}
\end{equation}
If only the first bound is known, one may set
$\epsilon_F=\sqrt{k}\epsilon_2$, where $k=\min(m,n)$.
Let $0\leq r<k$, and define
\begin{equation}
 L_r(Y)=\left(\sum_{j=r+1}^{k}
        (\sigma_j(Y)-\epsilon_2)_+^2\right)^{1/2}.
 \label{lr:lower}
\end{equation}
For any proposed globally defined behavior $b$, set
\begin{equation}
 B=H_b[S,E],\qquad U_b(Y)=\|Y-B\|_F+\epsilon_F.
 \label{lr:upper}
\end{equation}

\begin{theorem}[Evidence-certified representation-language advantage]
\label{lr:class-comparison}
On the event \eqref{lr:noise-event}, every real WFA $\mathcal A$ with
at most $r$ states satisfies
\[
 \|H_{f_{\mathcal A}}[S,E]-H_*\|_F\geq L_r(Y),
 \qquad \|B-H_*\|_F\leq U_b(Y).
\]
Consequently $U_b(Y)<L_r(Y)$ certifies that $b$ has strictly lower true
block squared loss than \emph{every} WFA with at most $r$ states.
The behavior $b$ may have been chosen using the same data $Y$.
\end{theorem}
\begin{proof}
Every such WFA block has rank at most $r$. Singular-value perturbation
gives $\sigma_j(H_*)\geq(\sigma_j(Y)-\epsilon_2)_+$, and the
Eckart--Young inequality bounds its distance from every rank-at-most-$r$
matrix below by $L_r(Y)$. The candidate upper bound is the triangle
inequality. Both statements hold simultaneously for all candidate
matrices on the fixed noise event, so selecting $b$ after seeing $Y$
does not invalidate this deterministic implication.
\end{proof}

\begin{remark}[What the certificate does not select]
The theorem does not discover $b$, compare the description lengths of
arbitrary languages, or prove that $r$ is a scientifically privileged
budget. It certifies a specific inequality between operationally
comparable behaviors. If $S,E$ or a stopping time are selected adaptively,
a fixed-block concentration theorem cannot simply be reused without
checking that its noise event remains valid. Confidence sequences,
simultaneous bounds, a valid union bound, or a fresh verification sample
are possible ways to discharge that separate statistical obligation.
\end{remark}

\begin{theorem}[Certified gated language replacement]
\label{lr:replacement}
Let $f_0$ be the behavior of an old real WFA with at most $r$ states, and
let $b$ be any globally defined presentation for which
$U_b(Y)<L_r(Y)$. Put
\[
 \Delta=L_r(Y)^2-U_b(Y)^2>0,
 \quad H_0=H_{f_0}[S,E],\quad D=B-H_0,
 \quad f_g=(1-g)f_0+gb.
\]
Then:
\begin{enumerate}
\item The family $f_g$ has a globally defined product-state realization
      for every real $g$. At $g=0$ it agrees with $f_0$ on every word.
\item For $\ell(H)=\tfrac12\|H-H_*\|_F^2$ and $0\leq g\leq1$,
\begin{equation}
 \ell(H_0+gD)-\ell(H_0)
 \leq-\frac{g\Delta}{2}
       -\frac{g(1-g)}{2}\|D\|_F^2.
 \label{lr:gate-bound}
\end{equation}
In particular every $0<g\leq1$ improves the true block loss, and
\[
 \left.\frac{d}{dg}\ell(H_0+gD)\right|_{g=0}
 \leq-\frac{\Delta+\|D\|_F^2}{2}<0.
\]
\item If the old presentation depends differentiably on parameters
      $\theta$, keeping the candidate fixed gives
      $\partial_\theta f_g=(1-g)\partial_\theta f_0$.
      Thus old derivatives are unchanged at initialization.
\item At $g=1$, projection to the candidate state preserves every
      transition and the observation map. The old subsystem may be
      discarded with no change to the resulting behavior on any word.
\end{enumerate}
\end{theorem}
\begin{proof}
Let $(X_0,x_0,T_a^0,o_0)$ and $(X_b,z_0,T_a^b,o_b)$ realize $f_0$ and
$b$. On $X_0\times X_b$, use the initial state $(x_0,z_0)$, transitions
$(x,z)\mapsto(T_a^0x,T_a^bz)$, and readout
$(x,z)\mapsto(1-g)o_0(x)+go_b(z)$. Induction on word length proves the
behavior formula. The candidate domain is forward invariant by the
global-presentation assumption.

For any two matrices the squared-norm identity gives
\[
 \ell((1-g)H_0+gB)
 =(1-g)\ell(H_0)+g\ell(B)
  -\tfrac12g(1-g)\|B-H_0\|_F^2.
\]
Theorem~\ref{lr:class-comparison} implies
$\ell(B)-\ell(H_0)\leq-\Delta/2$, proving
\eqref{lr:gate-bound}. Differentiating the exact identity gives the
stated gate derivative. The parameter identity follows by differentiation
with $b$ held fixed. Finally, at $g=1$ the readout ignores the old state
and the product transitions commute with projection onto $X_b$.
\end{proof}

\begin{corollary}[The initial descent direction is visible in measured loss]
\label{lr:observed-descent}
Under Theorem~\ref{lr:replacement}, put
$\widehat\ell(H)=\tfrac12\|H-Y\|_F^2$. Then
\[
 \left.\frac{d}{dg}\widehat\ell(H_0+gD)\right|_{g=0}
 \leq-\frac{\Delta+\|D\|_F^2}{2}<0.
\]
The inequality \eqref{lr:gate-bound} also holds with $\ell$ replaced by
$\widehat\ell$.
\end{corollary}
\begin{proof}
Since $H_0$ has rank at most $r$, the Eckart--Young inequality applied
to $Y$ gives
$\|H_0-Y\|_F\geq(\sum_{j>r}\sigma_j(Y)^2)^{1/2}\geq L_r(Y)$.
Also $\|B-Y\|_F\leq U_b(Y)$. Their squared-loss difference is at
most $-\Delta/2$. The same exact quadratic identity used in the
replacement theorem proves the assertions.
\end{proof}

\begin{remark}[The exact claim about new gradients]
The new gate is active. If $b$ also has trainable internal parameters,
their derivatives in this particular construction are multiplied by $g$
and vanish at $g=0$. This theorem does not claim simultaneous nonzero
gradients for every new parameter. The candidate may instead be fitted
before insertion, after which the gate supplies the initially active
direction. Nor does function preservation at $g=0$ mean that old
predictions remain fixed throughout the path: changing them is the
purpose of a successful revision. The stated line-loss bounds hold with
both branch behaviors fixed. Jointly training either branch changes that
path and requires a new argument or certificate.
\end{remark}

\begin{proposition}[Endpoint improvement is not path monotonicity]
\label{lr:monotonicity}
Under Theorem~\ref{lr:replacement}, $D\ne0$. The true unconstrained
line minimizer is
\[
 g_*={\langle H_*-H_0,D\rangle_F\over\|D\|_F^2}
 \geq{\Delta+\|D\|_F^2\over2\|D\|_F^2}.
\]
Consequently the true loss is nonincreasing on
\[
 0\leq g\leq g_{\rm safe}:=
 \min\left\{1,{\Delta+\|D\|_F^2\over2\|D\|_F^2}\right\}.
\]
It is nonincreasing on the entire interval $[0,1]$ if and only if
$\langle B-H_*,D\rangle_F\leq0$. A sufficient observed condition
for the latter is
\[
 \langle B-Y,D\rangle_F+\epsilon_F\|D\|_F\leq0.
\]
\end{proposition}
\begin{proof}
Strict candidate advantage excludes $B=H_0$. The line loss is a
strictly convex quadratic with derivative
$\langle H_0-H_*,D\rangle_F+g\|D\|_F^2$. Its initial derivative
bound from Theorem~\ref{lr:replacement} gives the stated lower bound
on $g_*$. The derivative is increasing, so its sign at $g=1$ gives
the necessary and sufficient condition. The observed sufficient
condition is Cauchy--Schwarz and the Frobenius noise bound.
\end{proof}

For example, scalar target $0$, old prediction $2$, and candidate $-1$
give squared error $(2-3g)^2$: every $0<g\leq1$ beats the old
endpoint, but the error rises after $g=2/3$. Thus a safe initial gate
step and full endpoint compression are distinct decisions. Stopping
before $g=1$ generally retains a genuine dependence on the old subsystem.

\begin{remark}[A small extension can cross a language boundary]
If the candidate has a one-dimensional rational state space, an
$r$-dimensional old WFA needs only one additional recurrent coordinate
during coexistence. The combined presentation can leave the linear-WFA
language through its transitions, its readout, or both. For the specific
$R_c$ candidate with $c>0$ and identity readout, the extra transition
coordinate is nonlinear. At the endpoint the realization reduces to the
candidate's one-dimensional state space. Thus the mechanism can decrease
the final state dimension while escaping a certified linear rank budget.
This does not compare total descriptive, computational, or numerical
complexity; those resources must be specified separately.
\end{remark}

\begin{corollary}[A nonvacuous certificate and replacement instance]
\label{lr:nonvacuous}
For any fixed $r\geq0$, choose $m>r$. Observe the behavior $b_1(a^n)=1/(n+1)$, with $x_0=1$, on all words
needed for this block. With exact evidence,
the candidate $b_1$ has $U_{b_1}=0<L_r(H_m)$. Therefore
Theorem~\ref{lr:replacement} applies to every old WFA with at most $r$
states. For fixed $m,r$, the same conclusion remains true for sufficiently
small noise whenever the candidate upper bound tends to zero.
\end{corollary}
\begin{proof}
Positive definiteness gives $\sigma_j(H_m)>0$ for all $j$, so
$L_r(H_m)>0$. The noisy assertion follows by continuity of singular
values and of the candidate upper bound at this strict inequality.
\end{proof}

\begin{corollary}[The harmonic mixture genuinely leaves the finite WFA language]
\label{lr:genuine-exit}
Let $b(a^n)=1/(n+1)$ and let $f_0$ have any finite WFA realization.
For every $g\ne0$, the behavior $f_g=(1-g)f_0+gb$ has no finite WFA
realization.
\end{corollary}
\begin{proof}
Finite WFA behaviors are closed under finite linear combinations by
block-diagonal direct sums and scalar changes of readout. If $f_g$ had
a finite realization, then $b=(f_g-(1-g)f_0)/g$ would as well, contrary
to Theorem~\ref{thm:cc-hilbert}.
\end{proof}

This structural conclusion concerns the specified candidate behavior.
The finite evidence certificate still establishes a target-relative
advantage only on its stated block and against its stated budget.

\begin{remark}[A real robustness limitation]
The qualification ``for fixed $m,r$'' matters. Hilbert matrices are
ill-conditioned, and higher exact Hankel rank is not automatically
observable above a fixed noise floor. The example does not establish a
uniformly robust certificate against arbitrarily large linear budgets.
Its strongest claim is a compact exact law plus budget-specific,
noise-qualified comparisons.
\end{remark}

\paragraph{A numerical certificate with its norm convention stated.}
Take $r=1$, $S=E=\{\varepsilon,a\}$, the measured block
\[
 Y=\begin{pmatrix}1&1/2\\1/2&1/3\end{pmatrix},
 \qquad B=Y,
\]
and suppose only that its unknown true block satisfies
$\|Y-H_*\|_F\leq\delta=10^{-3}$. This also gives the operator-norm
bound $\epsilon_2=\delta$. The singular values of $Y$ are
$(4\pm\sqrt{13})/6$, so
\[
 L_1={4-\sqrt{13}\over6}-10^{-3}
     \approx0.0647414540893,
 \qquad U_b=0.001,
\]
\[
 \Delta=L_1^2-U_b^2\approx0.00419045587760.
\]
Thus the candidate has strictly smaller true block loss than every
one-state WFA for every true block in the stated uncertainty ball. The
guaranteed half-squared-loss reduction at $g=1$ is at least
$\Delta/2\approx0.00209522793880$. If instead $10^{-3}$ is an
\emph{operator}-norm radius only, take $\epsilon_F=\sqrt2\,10^{-3}$;
then $U_b\approx0.00141421356237$ and
$\Delta\approx0.00418945587760$. No unobserved-word or physical-law
conclusion is part of this numerical certificate.

\section{The least observable extension and its precise scope}
\label{sec:localization}

The reciprocal coordinate also exposes the distinction between a small
state description and a small linear observable space. Work with rational
functions of $z$ on $(0,\infty)$ and the shift pullback
$\sigma(p)(z)=p(z+1)$. Let $S$ be the multiplicative set in $\mathbb R[z]$
generated by $z+n$ for $n\in\mathbb N$. Define
\[
 A=S^{-1}\mathbb R[z]
   =\mathbb R\bigl[z,(z+n)^{-1}:n\in\mathbb N\bigr].
\]
Every denominator is nonzero on the positive domain. The localization's
universal property is standard \cite{cc-stacks-localization}; the following
consequences give the exact extension required in this case.

\begin{theorem}[The least shift-stable reciprocal extension]
\label{thm:cc-localization}
The algebra $A$ is the least subalgebra of $\mathbb R(z)$ containing $z$
and $1/z$ and closed under $\sigma$. The shift extends uniquely to $A$.
More generally, for a commutative $\mathbb R$-algebra $B$ with $\mathbb R$-algebra endomorphism
$\tau$ and an element $b$ satisfying $\tau(b)=b+1$ and $b+n$ invertible
for all $n\ge0$, there is a unique $\mathbb R$-algebra homomorphism
$A\to B$ carrying $z$ to $b$ and commuting with the shifts.
The algebra $A$ is not finitely generated as a commutative
$\mathbb R$-algebra, although it is generated by $z$ and $1/z$ when the
shift operation is included in the signature.
\end{theorem}
\begin{proof}
Any shift-stable algebra containing $1/z$ contains
$\sigma^n(1/z)=1/(z+n)$, hence contains $A$. Conversely $\sigma(z)=z+1$
and $\sigma((z+n)^{-1})=(z+n+1)^{-1}$ lie in $A$, so $A$ is stable.
The universal property of localization extends $z\mapsto b$ uniquely,
because every member of $S$ maps to a unit. The extension commutes with
shift on $z$ and each inverse generator, hence on the algebra they generate.

For finite generation, every finite collection of rational functions has
poles in a finite union of finite sets. Polynomial combinations of those
functions cannot acquire a pole outside that union. But $A$ contains
$1/(z+N)$ for every nonnegative integer $N$, which has a pole at $-N$.
No finite set of algebra generators can produce them all. Finally, repeated
shift of the one generator $1/z$ produces every reciprocal generator, so
the enlarged operation signature has the claimed two-element generating set.
\end{proof}

The poles used in this proof lie outside the positive execution domain.
This is legitimate: equality of rational functions on the positive interval
implies equality as rational functions, so an identity on the execution
domain cannot manufacture a missing pole. No physical interpretation of
negative-state execution is required.

\paragraph{The specified observable and general reciprocal closure.}
The theorem adjoins the \emph{particular observable} $1/z$ and closes under
shift and algebra operations. It does not make $A$ closed under reciprocal
of every nowhere-zero observable. For example, $z+\tfrac12$ belongs to $A$
and is positive on the execution domain, but $(z+\tfrac12)^{-1}$ does not
belong to $A$: it has a pole at $-\tfrac12$, whereas members of $A$ can have
finite poles only at nonpositive integers. Thus $A$ is not the full rational
language used in Corollary~\ref{cor:cc-reciprocal}.

If closure under reciprocal of every function that is nowhere zero on the
\emph{whole positive domain} is required, use instead
\[
 P=\{p\in\mathbb R[z]:p(t)\ne0\text{ for every }t>0\},
 \qquad B=P^{-1}\mathbb R[z].
\]
This is the least algebra containing $\mathbb R[z]$ with that closure:
the requirement forces each $p\in P$ to be inverted, and a nowhere-zero
element $p/q\in B$ has $p\in P$, so its inverse $q/p$ again lies in $B$.
It is shift-stable because $p(z+1)\in P$ whenever $p\in P$.
Other certified execution domains give different localizations or partial
operations. In particular, an obligation holding only on one reachable
trajectory must not be silently strengthened to the whole positive domain.

This theorem makes a concrete point about changing available operations.
A countably generated ordinary observable algebra can have a finite
generating set once iteration of a specified operation is admitted. This
does not by itself assert a finite presentation by generators and relations. The new
language is economical in a different sense from a finite-dimensional linear
observable space. Neither economy is an invariant that can be identified
with all the others. For the running sequence, a one-dimensional
rational realization and infinite WFA dimension coexist without contradiction.

\section{Why state dimension alone cannot measure language cost}
\label{sec:cost}

\begin{proposition}[One smooth state can hide an arbitrary sequence]
\label{lr:smooth-one-state}
For every real sequence $(y_n)_{n\geq0}$ there is a smooth function
$g:\mathbb R\to\mathbb R$ such that the one-dimensional smooth
presentation
\[
 z_0=0,\qquad T(z)=z+1,\qquad o(z)=g(z)
\]
has unary behavior $f(a^n)=y_n$. If the sequence is bounded, $g$ can
be chosen so that each derivative of every fixed order is bounded on
$\mathbb R$.
\end{proposition}
\begin{proof}
Choose $\phi\in C_c^\infty((-1/3,1/3))$ with $\phi(0)=1$, and put
\[
 g(z)=\sum_{n=0}^{\infty}y_n\phi(z-n).
\]
The sum is locally finite and its summands have disjoint supports, so
it is smooth and $g(n)=y_n$. If $|y_n|\leq M$, then
$\|g^{(k)}\|_\infty\leq M\|\phi^{(k)}\|_\infty$ for every fixed
$k\geq0$.
\end{proof}

This construction stores the sequence in the readout. Its assertion is
existential and does not provide a finite description for an arbitrary
sequence. It shows why dimension and smoothness, even with bounded
readout derivatives for bounded sequences, are not by themselves a useful
cross-language complexity measure on this unbounded state domain.
A finite description language, numerical precision, computational cost,
and appropriate domain or regularity constraints are possible resources
to charge. Choosing them is part of the learning problem, not a
consequence of rank.

\section{Interaction with the whole task-specific elaborator}
\label{sec:elaboration}

A diagram of displayed categories can present the entire task-specific
elaborator: contextual input presentations, verification theories and their
models, syntax, telescopes, typed holes, admissibility obligations, evidence
provenance, interpretations, and proof-carrying transports. Such diagrams
need not be alternatives to that infrastructure; they can be its presentation.
Neither the scalar interface used here nor a chosen layering of displayed
structure is asserted to be a universal representation of all tasks. The companion note
\emph{A Task Specific Elaboration Situation for Changing Model Languages}
develops that wider presentation separately as a conditional specification
with typed records and telescopes, not an implemented full elaborator.
Displayed categories supply the modular categorical machinery
\cite{cc-displayed}; additional mathematical results must justify the
particular elaborator's existence, coherence, and any claimed canonicity.

This note contributes concrete nonvacuity tests to such a presentation.
For an input whose accepted law uniquely specifies $1/(n+1)$, the
finite-WFA compatible-description fiber is empty, while the rational fiber
contains the verified positive-domain recurrence. No section returning a
finite WFA on every such input can exist. Under only finite measurements,
the unrestricted WFA fiber is instead nonempty; a bounded-capacity
certificate can still rule out an explicitly budgeted class or prefer a
candidate on the tested block. The same categorical organization therefore
contains different successful and obstructed situations rather than
assuming that every hole can be filled.

Context restriction and language synthesis also have different roles.
Reindexing a verified description can be a cartesian lift when its syntax,
interpretation, and certificates reindex. Choosing a new characterization
that repairs an empty or inadequate fiber is an additional synthesis step;
it is not automatically a cartesian or opcartesian lift. The localization
of Section~\ref{sec:localization} is canonical only relative to its declared
operation, domain, and equations. It does not choose reciprocal as the
scientifically appropriate operation.

The interaction is testable in both directions. Verification constrains
which proposals count as globally defined candidates, which observations
are comparable, and which bounds may be used. A successfully elaborated
candidate then generates consequences such as the action-commutation test
of Proposition~\ref{lr:generated-tests}; a failed consequence can require a
new proposal. Canonicity of elaborated syntax, uniqueness of a represented
rational function, uniqueness of its coefficient chart, and empirical
identifiability of the target are distinct claims. Theorems above establish
only the particular claims whose premises they state.

\section{Close prior work and what remains unresolved}
\label{sec:prior}

Institution theory already supplies satisfaction-preserving vocabulary
change \cite{cc-goguen-burstall,cc-goguen-rosu}, and displayed categories
supply modular presentations of dependent categorical structure
\cite{cc-displayed}. Categorical automata learning and coalgebra learning
connect states, tests, and minimization \cite{lr:calf,lr:coalgebra-duality}.
The classical finite-Hankel-rank theory and its spectral learning methods
underlie the WFA obstruction and class comparison
\cite{cc-balle-mohri,lr:denis2016}. Localization provides the universal
algebraic construction \cite{cc-stacks-localization}. Braess and Hackbusch
already provide the substantially stronger exponential-sum approximation
rate discussed above \cite{cc-braess-hackbusch}.

Representation and law discovery are also established research subjects.
Symbolic regression searches equation forms as well as parameters
\cite{lr:schmidt2009}. SINDy selects sparse dynamics from a candidate
library and explicitly depends on coordinates and basis \cite{cc-sindy}.
Joint coordinate and equation discovery is addressed by Champion et al.
\cite{cc-sindy-coordinates}, and SINDy-PI addresses implicit dynamics and
rational nonlinearities \cite{cc-sindy-pi}. Crutchfield's hierarchical
machine reconstruction explicitly concerns innovation of computational
model classes \cite{cc-emergence}. These are relevant sources of candidate
proposals, not missing problems solved merely by redescribing them
categorically. Function-preserving architectural changes also have close
antecedents such as Net2Net \cite{lr:net2net}; the product-state mixture
used here is an elementary construction.

Two recent preprints are particularly close comparison targets.
Ma et al. jointly propose typed structural edits and diagnostic experiments,
include rational saturation and observation-map changes in their grammar,
and use anytime-valid evidence for class rejection \cite{lr:ma2026}.
Wang and Buehler develop verified regime transitions, evidence transport,
and proof-carrying provenance, including a Builder/Breaker example with a
description-length gate \cite{lr:wang2026}. Accordingly, this note does
not claim to introduce integrated representation revision, evidence
transport, or acceptance gates. Any priority claim even for the particular
sufficient-condition theorem would require a deeper literature audit.

What is proved here is a bounded synthesis: certified semantic-operation
nonclosure; positive-domain coordinate transport; exact, absolute-error,
and relative-error distinctions for the harmonic behavior; finite-data and
inverse-limit obstructions; conditional rational recovery with explicit
conditioning and domain obligations; a true-error class comparison that
licenses an executable gate and endpoint compression; and a precisely
scoped least observable extension. The common error and observation
interface makes these results depend on one another rather than merely
coexist in adjacent formalisms.

What remains unresolved is an effective and broadly useful procedure that
proposes an appropriate grammar, coordinate system, readout, or verification
theory from realistic evidence; compares their descriptive, numerical,
statistical, and computational costs; selects informative available tests;
and maintains the relevant certificates affordably. The finite-evidence
obstruction rules out an unrestricted finite-data decision that a failure
requires more width rather than a language change. The conditional positive
results license a specified change once its candidate and obligations are
available. Failure of a sufficient certificate does not prove that no useful
revision exists, and exact equality in a trace interface does not settle
all meanings relevant to a richer task.

\begin{thebibliography}{99}
\bibitem{cc-goguen-burstall}
J. A. Goguen and R. M. Burstall, Institutions: Abstract model theory for
specification and programming, \emph{Journal of the ACM} 39(1) (1992), 95--146.
\url{https://doi.org/10.1145/147508.147524}.

\bibitem{cc-goguen-rosu}
J. Goguen and G. Ro\c{s}u, Institution morphisms,
\emph{Formal Aspects of Computing} 13 (2002), 274--307.
\url{https://doi.org/10.1007/s001650200013}.

\bibitem{cc-displayed}
B. Ahrens and P. LeFanu Lumsdaine, Displayed categories,
\emph{Logical Methods in Computer Science} 15(1) (2019), article 20.
\url{https://lmcs.episciences.org/5252}.

\bibitem{cc-balle-mohri}
B. Balle and M. Mohri, Spectral learning of general weighted automata via
constrained matrix completion, \emph{Advances in Neural Information Processing
Systems} 25 (2012).
\url{https://papers.nips.cc/paper/2012/hash/700fdb2ba62d4554dc268c65add4b16e-Abstract.html}.

\bibitem{cc-stacks-localization}
The Stacks Project Authors, Localization, Section 10.9, especially
Proposition 10.9.3 (tag 00CM).
\url{https://stacks.math.columbia.edu/tag/00CM}.

\bibitem{cc-sindy}
S. L. Brunton, J. L. Proctor, and J. N. Kutz, Discovering governing equations
from data by sparse identification of nonlinear dynamical systems,
\emph{PNAS} 113(15) (2016), 3932--3937.
\url{https://doi.org/10.1073/pnas.1517384113}.

\bibitem{cc-sindy-coordinates}
K. Champion, B. Lusch, J. N. Kutz, and S. L. Brunton, Data-driven discovery
of coordinates and governing equations, \emph{PNAS} 116(45) (2019),
22445--22451.
\url{https://doi.org/10.1073/pnas.1906995116}.

\bibitem{cc-sindy-pi}
K. Kaheman, J. N. Kutz, and S. L. Brunton, SINDy-PI: A robust algorithm
for parallel implicit sparse identification of nonlinear dynamics,
\emph{Proceedings of the Royal Society A} 476 (2020), 20200279.
\url{https://doi.org/10.1098/rspa.2020.0279}.

\bibitem{cc-emergence}
J. P. Crutchfield, The calculi of emergence: Computation, dynamics and
induction, \emph{Physica D} 75 (1994), 11--54.
\url{https://csc.ucdavis.edu/~dynlearn/papers/CalcEmerg1_46.pdf}.

\bibitem{cc-braess-hackbusch}
D. Braess and W. Hackbusch, Approximation of $1/x$ by exponential sums in
$[1,\infty)$, \emph{IMA Journal of Numerical Analysis} 25(4) (2005), 685--697.
\url{https://doi.org/10.1093/imanum/dri015}.
Preprint: \url{https://www.mis.mpg.de/publications/preprint-repository/article/2004/issue-75}.


\bibitem{lr:denis2016}
F. Denis, M. Gybels, and A. Habrard, Dimension-free concentration bounds on
Hankel matrices for spectral learning, \emph{Journal of Machine Learning
Research} 17(31) (2016), 1--32.
\url{https://jmlr.org/papers/v17/14-501.html}.

\bibitem{lr:net2net}
T. Chen, I. Goodfellow, and J. Shlens, Net2Net: Accelerating learning via
knowledge transfer, arXiv:1511.05641 (2015).
\url{https://arxiv.org/abs/1511.05641}.

\bibitem{lr:schmidt2009}
M. Schmidt and H. Lipson, Distilling free-form natural laws from experimental
data, \emph{Science} 324 (2009), 81--85.
\url{https://doi.org/10.1126/science.1165893}.

\bibitem{lr:calf}
G. van Heerdt, M. Sammartino, and A. Silva, CALF: Categorical automata
learning framework, \emph{Computer Science Logic} (2017), article 29.
\url{https://doi.org/10.4230/LIPIcs.CSL.2017.29}.

\bibitem{lr:coalgebra-duality}
S. Barlocco, C. Kupke, and J. Rot, Coalgebra learning via duality,
\emph{Foundations of Software Science and Computation Structures} (2019),
62--79. \url{https://doi.org/10.1007/978-3-030-17127-8_4}.

\bibitem{lr:ma2026}
SiYuan Ma, Albert Gao, Chunzheng Zhu, Xin Yan, Wenlong Zhang, Wenxin Zhang,
Luqi Gong, Tianlin Li, and Qixin Zhang, Learning to outgrow a theory: Experimental discovery beyond
the initial hypothesis space, arXiv:2610.07627v1 (6 October 2026).
\url{https://arxiv.org/abs/2610.07627}.

\bibitem{lr:wang2026}
F. Y. Wang and M. J. Buehler, Self-revising discovery systems for science:
A categorical framework for agentic artificial intelligence,
arXiv:2606.01444v2 (revised 21 August 2026).
\url{https://arxiv.org/abs/2606.01444}.
\end{thebibliography}
\end{document}
